% Exact example 6: p=2, a=2, A=(16,8), B=(1/4,-1), P=(2,0).
\begin{tikzpicture}[font=\small,>=stealth]
\draw[rounded corners=3pt,gray!45] (0,0) rectangle (14.2,10.0);
\node at (7.1,9.45) {例 6：定点移到 $P(2,0)$，相同上侧正切对应下侧正切 $4$};
\begin{scope}[shift={(1.1,3.7)},scale=.5]
\draw[->,gray] (-.7,0)--(17.6,0) node[above] {$x$};
\draw[->,gray] (0,-4.1)--(0,9.0) node[left] {$y$};
\draw[teal!75!black,thick,samples=140,domain=-3.8:8.2,variable=\t]
 plot ({\t*\t/4},\t);
\draw[gray!75,dashed,thick] (.25,-1)--(16,8);
\draw[blue!75!black,very thick] (0,0)--(16,8);
\draw[orange!85!black,very thick] (0,0)--(.25,-1);
\draw[orange!85!black,densely dotted,thick] (.25,-1)--(.95,-3.8);
\draw[violet,thick] (2.8,0) arc[start angle=0,end angle=26.5650512,radius=2.8];
\node[violet] at (5.5,.9) {$\alpha$};
\draw[violet,thick] (3.9,0) arc[start angle=0,end angle=-75.9637565,radius=3.9];
\node[violet] at (4.0,-2.55) {$\beta$};
\fill (0,0) circle(2pt) node[above left=7pt] {$O$};
\fill (2,0) circle(2pt) node[below=8pt] {$P$};
\fill (16,8) circle(2pt) node[above left=7pt] {$A(16,8)$};
\fill (.25,-1) circle(2pt);
\node[anchor=east] at (-.5,-2.6) {$B$};
\draw[gray,thin] (-.45,-2.2)--(.20,-1.08);
\end{scope}
\node[anchor=west,align=left] at (10.0,6.3)
 {$y^2=4x$\\[12pt]$P(2,0)$\\[14pt]
 $\tan\alpha=\dfrac12$\\[15pt]
 $\tan\beta=\dfrac{2}{1/2}=4$};
\node at (7.1,1.05) {$B(\frac14,-1)$；灰色虚线 $AB$ 经过 $P$，两侧正切积为 $2$。};
\node at (7.1,.4) {与例 5 使用同一尺度；橙色点线为射线 $OB$ 的延长部分。};
\end{tikzpicture}
